\title{Large Language Models Can Automatically Engineer Features for Few-Shot Tabular Learning}

\begin{abstract}
Large Language Models (LLMs) have demonstrated outstanding capabilities in addressing complex and novel reasoning challenges, indicating their strong promise for tabular learning—a task fundamental to numerous practical scenarios. In this work, we introduce an innovative in-context learning strategy named \textsf{FeatLLM}, which leverages LLMs as feature engineering agents to generate enhanced input representations tailored for tabular prediction problems. The features synthesized by \textsf{FeatLLM} are subsequently input to straightforward downstream models, such as linear regression, enabling effective few-shot learning with strong predictive outcomes. Notably, our \textsf{FeatLLM} method requires only this basic predictive model and the constructed features during inference, avoiding repeated LLM calls for individual samples. Furthermore, it operates with simple API-level LLM access and addresses the issue of prompt length limitations. Extensive experiments spanning diverse tabular datasets illustrate that \textsf{FeatLLM} consistently produces high-quality, interpretable rules, achieving superior performance—on average, a 10\% improvement—over alternative frameworks like TabLLM and STUNT.
\end{abstract}

\section{Introduction}
Large Language Models (LLMs) have recently achieved notable success in producing coherent outputs for novel tasks, often with little to no fine-tuning specific to those tasks~\cite{creswell2022selection,imani2023mathprompter,taylor2022galactica}. Leveraging their broad exposure to vast text datasets, LLMs encapsulate a wealth of prior information that can serve as a stand-in for specialist knowledge in many domains~\cite{Genomes2021,singhal2023large,wilcox2003role}. This capacity has driven automation progress in varied fields such as dialogue analytics~\cite{finch2023leveraging} and automated code repair~\cite{fan2023automated}. By combining succinct task instructions and a few illustrative examples in prompt design, LLMs demonstrate the ability to grasp context and selectively emphasize critical features or samples~\cite{brown2020language,zhao2023survey}, underscoring their utility as advanced feature engineering tools.

Efforts to harness the rich reasoning and prior knowledge embedded in LLMs have extended into tabular learning applications. For example, medical predictions may benefit from the model’s understanding that age correlates with certain health risks. State-of-the-art techniques often express task objectives and feature definitions in natural language, translate data into serialized textual format, and then pass these prompts to an LLM~\cite{dinh2022lift,wang2023anypredict}. Subsequent model training is conducted through either in-context learning paradigms~\cite{nam2023semi} or efficient parameter tuning methods~\cite{dinh2022lift,hegselmann2023tabllm}. Incorporating LLMs’ embedded domain knowledge has been shown to yield gains over classical tabular models, especially in scenarios with limited labeled data~\cite{hegselmann2023tabllm}.

Existing methods that leverage LLMs for tabular data modeling face notable drawbacks. For approaches aiming at end-to-end prediction, each data instance necessitates at least one LLM inference, leading to significant computational overhead. Achieving high predictive accuracy generally requires fine-tuning the LLM, an impractical demand for models without publicly available training implementations or adaptation interfaces. This limitation is particularly pertinent as many leading LLMs recently released restrict interaction to inference APIs only~\cite{anil2023palm,openai2023gpt4}. Moreover, most of these strategies perform poorly when dealing with extensive prompts~\cite{liu2023lost}; as the number of tabular features grows and serialized input lengthens, these methods can become unusable or their efficacy substantially diminishes. Together, these issues fundamentally impede the adoption of LLM-based solutions for tabular learning in applied contexts.

To address these challenges, we reconceptualize the function of LLMs in the modeling pipeline: rather than relying on them exclusively for end-to-end inference, we employ LLMs for feature engineering, capitalizing on their embedded domain knowledge in a more computationally efficient fashion. Concretely, we propose an innovative in-context learning scheme, \textsf{FeatLLM}, which is designed to extract the "criteria" guiding LLM-based predictions. For example, in a medical prediction task, the LLM can articulate and identify explicit conditions—rules—that specify how feature patterns associate with the presence or absence of a disease. These inferred criteria, in turn, enable the construction of new, informative features that serve as alternatives to the original feature set for downstream modeling. This paradigm shifts the LLM’s utility toward generating higher-level representations, thereby facilitating the use of much simpler downstream models. Once meaningful features are derived, resource-intensive inference with sophisticated models is no longer necessary, thereby reducing inference time. The capacity of LLMs for in-context learning allows feature discovery by incorporating labeled examples into prompts—without necessitating additional training. Additionally, methods borrowed from ensemble learning~\cite{breiman2001random}, such as feature bagging, offer practical solutions to manage prompt size when the dimensionality of the tabular data is large.

\textsf{FeatLLM} organizes the process of engineering new features through input prompt design into two main stages: (1) interpreting the given problem and deducing the links between input features and target labels; and (2) using the insights from the initial stage, along with several example instances, to synthesize explicit rules that can effectively differentiate between classes. This stepwise reasoning strategy enables LLMs to disregard features that do not contribute meaningful information during rule formulation. These formulated rules are subsequently applied to the dataset—by parsing them as executable code—which allows for the transformation of original data into binary indicators representing whether each sample satisfies a given rule. A simple, low-complexity classifier is then trained on these binary features to predict class membership probabilities. To bolster robustness, an ensemble methodology is adopted where feature generation is iteratively conducted with feature bagging. By controlling both the quantity of features and samples selected during bagging, the limitation imposed by large prompt sizes is alleviated. \textsf{FeatLLM} operates without the need for additional model training and incurs minimal inference overhead, thereby making LLM-based solutions applicable across a variety of practical settings.

Our empirical assessment of \textsf{FeatLLM} spans 13 tabular datasets under low-shot conditions, demonstrating both consistently strong and reliable results. The findings illustrate that LLMs are capable of leveraging pre-existing as well as data-driven knowledge to derive informative features. Our approach surpasses leading few-shot learning baselines across multiple scenarios. The codebase has been made accessible via an anonymized GitHub repository at~\texttt{https://github.com/Sungwon-Han/FeatLLM}.

\section{Related Work}

\subsection{Few-Shot Learning with Tabular Data} Developments in few-shot learning have made it possible to derive robust, generalizable data representations with only a handful of labeled examples, thereby minimizing labeling expenditure~\cite{chen2018closer}. While this paradigm initially concentrated on the image domain, its high capacity for generalization has now been evidenced in diverse modalities, including text, audio, and more recently, tabular datasets~\cite{majumder2022few,nam2023stunt,schick2021exploiting}. Training with very limited labeled data, however, tends to heighten susceptibility to spurious pattern learning~\cite{bartlett2021deep}. To address these issues, recent methods have exploited auxiliary data sources or incorporated priors tailored to the application domain, effectively embedding useful inductive biases within the training process. Approaches include leveraging vast amounts of unlabeled tabular data with well-designed augmentation methods for semi-supervised training~\cite{ucar2021subtab,yoon2020vime}, and harnessing generic unsupervised pre-training to achieve superior model initialization~\cite{somepalli2021saint}. Unsupervised learning objectives often center around masking or corrupting inputs followed by reconstruction tasks~\cite{arik2021tabnet,majmundar2022met}, or around utilizing data augmentation in combination with contrastive loss frameworks~\cite{bahri2022scarf,chen2020simple}. Yet, because tabular data frequently encompasses heterogeneous features, devising augmentation strategies that universally benefit learning remains a substantial challenge, with most augmentation-based approaches showing limited improvement for few-shot tasks~\cite{nam2023stunt}.

More recently, researchers have investigated the effectiveness of pre-training models on diverse tabular datasets, not limited to the target task domain, followed by transfer to the target application~\cite{zhang2023generative,levin2022transfer,zhu2023xtab}. As an illustrative example, TransTab~\cite{wang2022transtab} employs transformer-based architectures to learn semantic interactions among table columns through the use of column names rather than raw table contents. This strategy allows the accumulated knowledge from disparate tables and fields to be applied to unseen tabular tasks, achieving zero-shot and few-shot learning performance. In parallel, TabPFN~\cite{hollmann2022tabpfn} synthesizes tabular datasets by sampling from multiple distribution types, utilizing these for the pre-training of a Bayesian neural network. Upon completion of pre-training, both training and test data from the target problem are supplied for inference directly, obviating the need for additional training. The transfer of prior information distilled from many source datasets has proven to outperform classic tree-based algorithms and demonstrates significant efficacy under few-shot regimes.

\subsection{Language-Interfaced Tabular Learning} Integrating large language models (LLMs)~\cite{anil2023palm,openai2023gpt4} presents another valuable avenue for importing prior knowledge into tabular learning tasks. Trained on extensive textual datasets encompassing broad linguistic and conceptual information, LLMs have established their ability to perform a wide array of novel tasks via transfer~\cite{brown2020language}. Recent advancements focus on utilizing this knowledge-rich foundation for tabular data analysis by converting tabular inputs into textual representations suitable for LLM prompt formats~\cite{dinh2022lift,hegselmann2023tabllm,wang2023anypredict}. By supplying the model with encoded tables, descriptive metadata, and contextual task statements as prompts, these methods enable inference by aligning feature relevance and task objectives through the LLM’s reasoning ability. Strategies for leveraging LLMs further include model adaptation through lightweight parameter-efficient fine-tuning methods such as LoRA~\cite{hu2021lora} and IA3~\cite{liu2022few}, as well as utilizing in-context learning, which incorporates a limited number of training examples within the input prompt to guide predictions, all without updating model weights~\cite{dinh2022lift,hegselmann2023tabllm,nam2023semi,slack2023tablet}.

Although the previously discussed approaches enhance performance in few-shot tabular learning, their reliance on executing inference through the LLM for every prediction introduces practical obstacles for deployment within industry settings. In contrast, this research seeks to move beyond the traditional end-to-end black-box paradigm, in which an LLM processes each data point directly for classification. The proposed method instead harnesses the LLM to infer class-specific rules or conditions. \textsf{FeatLLM} subsequently interprets these derived rules into executable code, constructing novel features that facilitate downstream predictions independently of the LLM. This approach employs in-context learning to extract rules by leveraging existing knowledge in conjunction with limited labeled samples.

\section{Methods}
The central aim of \textsf{FeatLLM} is to automatically discern ``rules"—that is, the sets of criteria linked to each output class—from a small collection of labeled examples, as illustrated in Figure~\ref{fig:main_model}. Rules are initially inferred by the LLM, after which they are converted into programmatic statements to yield new binary indicators for every input, each signifying whether a particular rule holds true. These binary variables collectively serve as input to a weighted sum (dot product) with non-negative, trainable parameters, providing a measure of class probabilities. Through model training, one learns the relative significance of each rule-feature in predicting the corresponding classes (Section~\ref{Sec:training}). This methodology is iteratively applied with bagging and the outputs are integrated via ensemble techniques. Details on the formal problem setup and the individual methodological components are provided below.

\vspace*{-0.12in}
\paragraph{Problem formulation.} Consider a tabular dataset comprising $N$ labeled examples, represented as $\mathcal{D} = \{(\mathbf{x}^i, \mathbf{y}^i)\}_{i=1}^{N}$. Each data point $\mathbf{x}^i$ is a $d$-dimensional vector corresponding to $d$ distinct input features, and the feature names—articulated in natural language—are given by $F = \{f_j\}_{j=1}^{d}$. Accompanying these feature names are their respective definitions, supplied as supplementary information. In a classification context, the label $\mathbf{y}^i$ belongs to one of several predefined classes. In the $k$-shot learning scenario, we select a subset of $k$ ($k < N$) labeled examples at random from $\mathcal{D}$ to train our model.

\subsection{Prompt Design for Extracting Rules}
\label{Sec:prompt_design}
To foster more precise reasoning when extracting rules with the LLM, we shape the task to parallel the analytical process of a domain expert addressing tabular learning problems. Our prompt structure consists of the following three core elements (illustrated in Fig.~\ref{fig:prompt_for_rule} and detailed in Appendix~\ref{sec:example_rule}):

\vspace*{-0.12in}
\paragraph{Basic information description.} This section delivers all pertinent details necessary to resolve the task at hand (highlighted in orange in Fig.~\ref{fig:prompt_for_rule}). It incorporates a statement of the task—alongside label information—clear definitions for each feature, and illustrative sample instances. The task statement is typically posed as a direct question (for example, ``Does this patient's myocardial infarction complications data suggest chronic heart failure? Yes or no?"). Additional examples of such formulations appear in Appendix~\ref{sec:task_desc}. For feature descriptions, both the datatype and, where available, a concise definition are included. When features are categorical, representative categories may be specified as examples. To ground the model, we serialize a limited number of training samples, pairing them with their actual labels. The serialization format aligns with conventions from previous literature~\cite{dinh2022lift,hegselmann2023tabllm}:

\begin{align}
&\text{Serialize}(\mathbf{x}^i,\mathbf{y}^i, F) = \nonumber \\ 
&``\ f_1\ \text{is}\ \mathbf{x}^i_1.\ ... \ f_d\  \text{is}\  \mathbf{x}^i_d.\ \ \text{Answer:}\  \mathbf{y}^i\ ”
\end{align}

\vspace*{-0.12in}
\paragraph{Reasoning instruction.} Our objective is to improve the reasoning abilities of the LLM by supplying targeted reasoning instructions (refer to the blue text in Fig.~\ref{fig:prompt_for_rule}). The instruction opens with a statement akin to the chain-of-thought methodology~\cite{wei2022chain}. We structure the inference process undertaken by the LLM into two distinct phases. During the initial phase, the model is prompted to deduce the potential causal associations or trends between the provided features and the task, independent of any exemplar demonstrations and based solely on domain knowledge or commonsense reasoning. This promotes comprehensive utilization of the LLM's inherent prior knowledge regarding the task. In the subsequent phase, the LLM integrates illustrative demonstrations alongside insights from the first phase to generate class-specific rules. Adopting this bifurcated reasoning sequence minimizes the risk of the model latching onto spurious patterns present in extraneous columns and steers its attention toward more pertinent features. To regulate the comprehensiveness of rule extraction and avert issues of underfitting or an unwieldy prompt size, we stipulate a set quantity of rules (namely, 10)\footnote{Further discussion on rule count appears in Section~\ref{sec:analysis}.} to be derived for each class. Furthermore, when developing these rules, the LLM is directed to consult the feature type as specified in the basic information section to preclude any mismatches in rule-type correspondence.

\vspace*{-0.12in}
\paragraph{Response instruction.} To streamline subsequent parsing and application of the extracted rules, we explicitly instruct the LLM on the organization of its responses (see the yellow-highlighted content in Fig.~\ref{fig:prompt_for_rule}). The prompt articulates guidelines on the designated response formats for each answer category (i.e., response formatting) and clarifies the prescribed template for structuring individual rules (i.e., rule formatting). Through these directives, the LLM can adapt its response style to align with the types of features present. In the absence of such specific guidance, the LLM might employ logical connectors such as AND or OR to combine multiple rules. An illustrative sample of the generated output can be found in Appendix~\ref{sec:outcome-rule}.

\subsection{Inferring Class Likelihood via Rules}
\label{Sec:training}

\paragraph{Parsing rules for feature generation.} The rules obtained from the earlier phase are employed to construct new binary features. Specifically, for each class, a feature is generated to indicate whether a given sample satisfies the class-specific rules (taking a value of either '0' or '1'). These features assist in estimating the likelihood that a sample falls into a particular class. Since the LLM produces rules in natural language, an automated parsing mechanism must be devised for effective feature extraction. Although we standardize the rule formats during generation to facilitate parsing, the LLM's outputs can still exhibit unwanted syntactic variations~\cite{kovavc2023large}. For example, the LLM may replace symbols such as ``$>$” or ``$<$” with their textual equivalents, such as ``greater than'' or ``less than.'' These inconsistencies further complicate the parsing process.

To overcome the difficulties of parsing syntactically inconsistent outputs, we opt against writing extensive parsing programs and instead utilize the LLM itself for this purpose. Owing to its proficiency in understanding the semantics of text amid varied syntax and its capacity to generate code from textual instructions (benefiting from code-centric training data), the LLM can facilitate this step. We prepare a prompt that contains the function name, explanations for both the input and output, and the inferred rules, which is then provided to the LLM. Illustrative examples of such prompts and resulting outputs can be found in Figures~\ref{fig:prompt_for_parsing} and ~\ref{sec:example_parsing}-\ref{sec:example_parse_outcome} in the Appendix. The code produced by the LLM is subsequently executed using Python’s \texttt{exec()} function to convert the data.

\vspace*{-0.12in}
\paragraph{Inferring class likelihood.} Once binary features corresponding to class-specific rules have been generated, an intuitive approach for estimating the likelihood that a sample belongs to each class is to count the number of class rules satisfied by the sample (i.e., summing the binary features per class). Nevertheless, since the informativeness of rules and features can differ, it is necessary to quantify their relative importance using training data. We address this by fitting a standard machine learning (ML) model to the binary feature representation for each class. As an example, a linear model with no bias term may be employed. Let $\mathbf{z}_k^i \in \{0, 1\}^R$ represent the binary feature vector generated for sample $\mathbf{x}^i$ with respect to class $k$, where $R$ denotes the count of rules per class and $c$ signifies the number of classes. The resulting class probabilities $\mathbf{p}^i$ are computed as follows:

\begin{align}
\text{logit}_k^i &= \max(\mathbf{w}_k, 0) \cdot \mathbf{z}_k^i, \nonumber \\
\mathbf{p}^i &= \text{Softmax}({[\text{logit}_{1}^i, ..., \text{logit}_{c}^i]}). \label{eq:infer_prob}
\end{align}

In this framework, $\mathbf{w}_k$ denotes the trainable parameters in the linear model, encapsulating the contribution of each input feature. By constraining these weights within a positive subspace, we guarantee that LLM-derived features do not decrease the class probability during training. The output logit vector is subsequently transformed into class probabilities using the softmax operator. The parameters $\mathbf{w}_k$ are then optimized by minimizing the cross-entropy loss based on a limited number of few-shot training examples.

\vspace*{-0.12in}
\paragraph{Ensembling with bagging.} The methodology involves repeating the above process multiple times to form several inference models, which are then aggregated through ensembling to determine the final class prediction. After $T$ trials, we amass a set of learned weights, $\{\mathbf{w}[t]\}_{t=1}^T$, and compute the predicted class probability by averaging individual predictions $\mathbf{p}^i$ produced using each set of weights $\mathbf{w}[t]$, as formalized in Eq.~\ref{eq:infer_prob}.

To encourage diversity in rule extraction for the ensemble, we employ an elevated temperature setting during LLM inference. Additionally, the sequence of few-shot demonstration examples is randomized, since prior research indicates that their order can influence the responses of the LLM~\cite{lu2022fantastically}\footnote{Further analysis of rule diversity appears in Appendix~\ref{sec:diversity}}. To harness this diversity for greater predictive reliability, we use bagging—randomly selecting subsets of features or data instances for each trial—which is particularly advantageous when the dataset includes numerous training examples or features. The ensembling strategy offers two key strengths. Firstly, even when certain runs yield erroneous rules from the LLM, accurate rules identified in other runs can counterbalance these errors, bolstering framework robustness. This is aligned with the LLM's property of self-consistency~\cite{wang2022self}, whereby rules emerging frequently across separate trials tend to be reliable. Secondly, ensembling directly mitigates the prompt size constraint inherent to LLMs: with tabular data that might exceed input length limits, bagging enables reduction of input size, maintaining strong performance. While an individual rule set may miss interactions among all features, the ensemble—which combines several weak learners produced over multiple runs—ameliorates shortcomings due to incomplete feature or instance representation in any single trial.

\section{Experiments}
We assess the effectiveness of \textsf{FeatLLM} across a variety of tabular datasets and carry out a series of studies to better understand the operation of our proposed framework. For brevity, supplementary experiments—including tests on different LLMs, qualitative assessments, and other evaluations—are included in the Appendix.

\subsection{Performance Evaluation}
\paragraph{Datasets.}
Our experimental setup employs 13 tabular datasets tailored for binary or multi-class classification challenges: (1) Adult~\cite{asuncion2007uci}, which involves forecasting whether an individual's annual income surpasses \$50,000; (2) Bank~\cite{moro2014data}, targeting the likelihood that a customer subscribes to a term deposit; (3) Blood~\cite{yeh2009knowledge}, focused on predicting donor return behavior; (4) Car~\cite{kadra2021well}, where the goal is to classify the quality rating of vehicles; (5) Communities~\cite{misc_communities_and_crime_183}, used for estimating crime rates within geographic areas; (6) Credit-g~\cite{kadra2021well}, predicting creditworthiness; (7) Diabetes\footnote{\texttt{kaggle.com/datasets/uciml/pima-indians-diabetes-database}}, which forecasts diabetes diagnoses; (8) Heart\footnote{\texttt{kaggle.com/datasets/fedesoriano/heart-failure-prediction}}, for predicting coronary artery disease occurrence; and (9) Myocardial~\cite{misc_myocardial_infarction_complications_579}, which determines chronic heart failure presence.

Additionally, our experiments incorporate more contemporary and synthetic datasets that have received limited attention in prior studies and were not part of the training data for LLMs: (10) Cultivars, aiming to predict the potential yield of soybean cultivars\footnote{\texttt{archive.ics.uci.edu/dataset/913}}; (11) NHANES, designed for classifying individuals' age groups from recorded features\footnote{\texttt{archive.ics.uci.edu/dataset/887}}; (12) Sequence-type, which classifies numeric sequences as arithmetic, geometric, Fibonacci, or Collatz types; (13) Solution-mix, where the task is to predict if a mixture's percentage concentration from four components exceeds 0.5. The Cultivars and NHANES datasets were added to the UCI Machine Learning Repository after September 2023, confirming their exclusion from the pre-training data of the LLM API (gpt-3.5-turbo-0613) utilized in our experiments. The latter two datasets are artificially generated to prevent any potential memorization by the LLM and to rigorously test evaluation.

The datasets in use present a broad spectrum of size and complexity, as outlined in Table~\ref{Tab:basic-desc}. Each dataset offers explicit attribute names and descriptions. Some datasets, such as `Communities' and `Myocardial', comprise more than 100 features, which creates additional difficulties given the context window constraints typical for LLMs.

\vspace*{-0.12in}
\paragraph{Baselines.}
Our empirical comparisons include nine baseline models. The initial group of three baselines consists of standard methods for tabular data: (1) Logistic regression (LogReg), (2) XGBoost~\cite{chen2016xgboost}, and (3) RandomForest~\cite{ho1995random}. Two further baselines assume the availability of large unlabeled datasets: (4) SCARF~\cite{bahri2022scarf} and (5) STUNT~\cite{nam2023stunt}. It is important to acknowledge that acquiring extensive unlabeled data may be impractical for many real-world scenarios. We also compare against (6) TabPFN~\cite{hollmann2022tabpfn}, which relies on pretraining via synthetically generated data with varying distributions. The last three baselines leverage LLM-based methodologies: (7) In-context learning (In-context)~\cite{wei2022emergent}, (8) TABLET~\cite{slack2023tablet}, and (9) TabLLM~\cite{hegselmann2023tabllm}. In-context learning integrates few-shot training examples into the prompt directly, requiring no parameter tuning. TABLET augments prompts with supplementary information provided by an external classifier through rule-sets and prototype examples to facilitate more robust predictions. TabLLM adopts a parameter-efficient strategy (such as IA3~\cite{liu2022few}) for tuning the LLM using few-shot data.

\vspace*{-0.12in}
\paragraph{Implementation details.}
\textsf{FeatLLM} is implemented with GPT-3.5 as the primary LLM backbone but remains flexible with respect to alternative LLM backbones (refer to Appendix~\ref{sec:backbones} for experimental outcomes using various backbones). For inference, we use a temperature setting of 0.5 and keep the top-p parameter at the API default value of 1. The ensemble count and rule extraction count are configured to 20 and 10, respectively. Further discussion of hyper-parameter effects appears in Figure~\ref{fig:hyperparameter2} and Appendix~\ref{sec:hyper-parameter}. The linear model training utilizes the Adam optimizer with a learning rate of 0.01 and is run for 200 epochs. To select the optimal epoch, we perform $k$-fold cross-validation. For baseline comparisons, in-context learning employs GPT-3.5, while TabLLM follows the original configuration by employing T0~\cite{sanh2022multitask}. Note that due to TabLLM’s limitation to LLMs with publicly released checkpoints, GPT-3.5 is not supported in these experiments. For traditional machine learning approaches—including LogReg, XGBoost, and RandomForest—we use Grid-search along with $k$-fold cross-validation to select optimal parameters. The $k$ value is chosen as either 2 or 4, ensuring each training set contains examples from every class. Additional implementation specifics can be found in Appendix~\ref{sec:baselines}.

\vspace*{-0.12in}
\paragraph{Results.}
Tables~\ref{tab:main1} and \ref{tab:main2} display comparative performance outcomes across multiple datasets. Each experiment was conducted three times, varying the random splits of the training and test sets; the mean and standard deviation of the resulting AUC scores are reported. Across all comparisons, \textsf{FeatLLM} either leads or is among the top two performing approaches, surpassing most baselines. Remarkably, our method achieves similar levels of performance with only 4 shots as traditional tabular learning baselines do with 32 shots, as illustrated in Fig.~\ref{fig:compares}. Although both STUNT and TabPFN yield substantial gains for particular datasets, their aggregate improvement over classic machine learning techniques is relatively limited. Furthermore, approaches that utilize LLMs—such as in-context learning, TABLET, or TabLLM—were unable to successfully process tabular datasets characterized by high feature dimensionality. By integrating bagging and ensembling, our framework effectively addresses datasets with many features and achieves competitive results. It should also be highlighted that this bagging and ensemble methodology could be incorporated into preexisting LLM-based models. Nevertheless, such integration would double inference operations per instance, greatly increasing computational demands and rendering the approach less feasible in practice. 

\subsection{Analysis \& Discussion}
\label{sec:analysis}
\paragraph{Ablation study.} To assess the contributions of key framework components, we conducted ablation experiments examining: (1) \textsf{-Tuning}: the absence of the linear model's weight tuning, instead relying on summing binary features to calculate the logit; (2) \textsf{-Ensemble}: removal of the ensembling step; (3) \textsf{-Description}: excluding feature descriptions; and (4) \textsf{-Reasoning}: removing Step-1 within the reasoning instruction when deriving rules.

Table~\ref{tab:ablation} presents the dataset-averaged changes in AUC following each ablation. In every scenario, the exclusion or modification of a key component resulted in diminished performance. Notably, the contribution of weight tuning to performance improvement becomes more significant as the shot count rises, implying that larger sample sizes facilitate more precise rule weighting. In contrast, both feature descriptions and explicit reasoning instructions yield greater benefits when training data is scarce, emphasizing the role of leveraging LLM’s prior knowledge in few-shot settings. Finally, across all configurations, the proposed ensemble technique consistently yields enhanced performance.

\vspace*{-0.12in}
\paragraph{Training \& inference time comparison.} We evaluate and contrast the time required for both training and inference across various approaches. All experiments were carried out on the Adult dataset with 16-shot training sets. A single A100 GPU serves as the default computational resource, apart from TabLLM, which employs four A100 GPUs to enable model parallelism. The corresponding runtime measurements are presented in Table~\ref{tab:cost}. A key observation is that \textsf{FeatLLM} achieves inference times on par with traditional baselines, maintaining relatively efficient performance. In stark contrast, alternative LLM-oriented techniques—including In-context learning, TABLET, and TabLLM—exhibit significantly greater inference times. Regarding training time, \textsf{FeatLLM} is competitive with other LLM-based methods: it necessitates only 30 API calls, representing a modest computational expense that is easily parallelizable.

\vspace*{-0.12in}
\paragraph{Quality of features generated by \textsf{FeatLLM}.}
We design a straightforward test to assess how informative the rule-based features produced by \textsf{FeatLLM} are for real-world tasks. Our evaluation quantifies the average AUROC score between the newly constructed features and the corresponding target labels. Subsequently, we determine the proportion of features deemed useful, characterized by an AUROC exceeding 0.5 for each dataset. 
A summary of these results is reported in Table~\ref{tab:quality}. The findings reveal that a significant fraction of the generated rules are indeed pertinent to the task-specific target, thereby offering substantial information for the main problem (see the “Before tuning” column). Additionally, once the linear model undergoes tuning with empirical data, rules that exhibit minimal data relevance (defined as those assigned weights below a given threshold) are automatically disregarded (see the “After tuning” column). The threshold was established at 0.1, based on the global weight distribution histogram.

\vspace*{-0.12in}
\paragraph{Effect of adding spurious correlations.} To evaluate the ability of different approaches to mitigate the influence of extraneous spurious correlations under low-shot conditions, we performed a targeted investigation. Our setup relied on the Heart dataset, artificially augmenting it by appending randomly selected, task-irrelevant columns sourced from the Adult dataset. We assessed how model performance varied as the quantity of these non-informative features increased. In order to ensure an equitable assessment, each method was granted access to an identical 8-shot set of labeled data, deliberately omitting the use of unlabeled samples and thus excluding techniques such as SCARF and STUNT.

Figure~\ref{fig:spurious} depicts the impact of spurious features on predictive accuracy. Relative to competing methods, \textsf{FeatLLM} exhibited the smallest decline in performance as more irrelevant columns were included. This resilience can likely be attributed to our framework’s mechanism, which incorporates prior domain knowledge to restrict rule generation to features demonstrably related to the main task—consequently minimizing the influence of irrelevant features and bolstering robustness. Empirically, we found that rules associated with noisy columns represented just 1–2\% of all rules extracted. Nonetheless, our analysis also reveals that if columns from other semantically similar datasets are appended, \textsf{FeatLLM} may occasionally internalize spurious patterns and misconstrue the causal links between the features and the task (see Appendix~\ref{sec:spurious-appendix}).

\vspace*{-0.12in}
\paragraph{Hyper-parameter analysis}
\textsf{FeatLLM} relies primarily on two hyper-parameters: the number of ensemble members and the number of rules per class generated in each LLM interaction. To examine how these hyper-parameters affect overall model performance, we performed a series of experiments. Our observations indicate that increasing the ensemble size generally enhances performance up to a threshold—approximately 20 ensembles—after which improvements plateau (see Fig.~\ref{fig:appendix-hyperparameter1} in the Appendix). Regarding the count of rules per class, statistical analysis revealed no significant variations in performance; nevertheless, selecting 10 rules provided the optimal balance between the length of the prompt and model accuracy (see Fig.~\ref{fig:hyperparameter2}). We hypothesize that augmenting the number of rules increases the dimensionality of inputs, offering more information but potentially leading to overfitting, especially when data is scarce.

\section{Conclusion}
We introduce \textsf{FeatLLM}, a method that pushes the boundaries of LLM-driven few-shot learning for tabular data. Rather than limiting LLMs to direct prediction, \textsf{FeatLLM} leverages their capabilities to create interpretive rules—akin to automated feature engineering—that guide downstream predictions. By incorporating ensemble strategies such as bagging and rule generation, our approach operates efficiently with minimal inference costs, requires no additional training (relying solely on API access), and sidesteps constraints related to the number of input features. The resulting model demonstrates strong predictive accuracy and constitutes an effective and resource-efficient strategy for applying LLMs to real-world tabular data.

Currently, \textsf{FeatLLM} is tailored for situations involving limited labeled data (low-shot scenarios). Moving forward, our aim is to extend its adaptability for larger datasets, integrate more diverse feature construction mechanisms beyond rule extraction, and improve interpretability—broadening its suitability across more task domains.

\section{Broader Impact}
The architecture of \textsf{FeatLLM} is crafted to exploit LLMs’ prior knowledge, with the goal of surpassing what is typically attainable through conventional supervised approaches for tabular learning. By effectively harnessing prior information, our research addresses the key challenge of overfitting when data is scarce and enhances data-efficient learning. \textsf{FeatLLM} is especially valuable in practical settings such as finance and healthcare, where obtaining labeled data is labor-intensive or costly, and where LLMs’ broad pretraining (often covering finance and healthcare texts) becomes a significant asset. A critical avenue for future work involves evaluating and mitigating potential societal biases or misinformation encoded in LLMs’ prior knowledge, as this external information can heavily influence, or even overshadow, learning from observed labels.

\end{document}